% Preámbulo común del material docente · XeLaTeX
\documentclass[11pt,a4paper]{article}
\usepackage[margin=22mm,top=20mm,bottom=22mm]{geometry}
\usepackage{fontspec}
\setmainfont{Avenir Next}[BoldFont={Avenir Next Demi Bold}, BoldItalicFont={Avenir Next Demi Bold Italic}]
\setsansfont{Avenir Next}
\setmonofont{Menlo}[Scale=0.88]
\newfontfamily\symfont{Apple Symbols}
\newfontfamily\unifont{Arial Unicode MS}
\usepackage{unicode-math}
\setmathfont{latinmodern-math.otf}
\usepackage{polyglossia}\setdefaultlanguage{spanish}
\usepackage{xcolor,booktabs,tabularx,enumitem,parskip,microtype,titlesec,fancyhdr}
\usepackage[most]{tcolorbox}
\usepackage[hidelinks]{hyperref}
\emergencystretch=2em

\definecolor{tinta}{HTML}{1F2933}
\definecolor{acento}{HTML}{B7410E}
\definecolor{suave}{HTML}{F4F1EC}
\definecolor{gris}{HTML}{6B7280}
\definecolor{linea}{HTML}{D6D0C6}
\color{tinta}

\titleformat{\section}{\Large\bfseries\color{acento}}{\thesection}{0.6em}{}
\titleformat{\subsection}{\large\bfseries}{\thesubsection}{0.6em}{}
\titlespacing*{\section}{0pt}{1.6em}{0.5em}
\titlespacing*{\subsection}{0pt}{1.1em}{0.3em}
\setlist{nosep, leftmargin=1.4em, itemsep=2pt}
\renewcommand{\arraystretch}{1.25}

% bloque de órdenes de terminal
\newtcblisting{orden}{listing only, colback=suave, colframe=linea, boxrule=0.4pt, arc=2pt,
  left=6pt, right=6pt, top=4pt, bottom=4pt, listing options={basicstyle=\ttfamily\small, breaklines=true, columns=fullflexible, keepspaces=true}}
% nota destacada
\newtcolorbox{nota}[1][Nota]{colback=white, colframe=acento, boxrule=0.6pt, arc=2pt, title=#1,
  fonttitle=\bfseries\small, coltitle=white, colbacktitle=acento, left=6pt, right=6pt}
% preguntas
\newcommand{\pregunta}[2]{\item[\textbf{(#1)}] #2}
\newenvironment{preguntas}{\begin{description}[leftmargin=2.4em, style=nextline, itemsep=4pt]}{\end{description}}
\newcommand{\ok}{{\symfont ✔}}
\newcommand{\nok}{{\symfont ✘}}
\newcommand{\qed}{{\symfont ∎}}
\newcommand{\codigo}[1]{\texttt{#1}}

\pagestyle{fancy}\fancyhf{}
\renewcommand{\headrulewidth}{0pt}
\fancyfoot[C]{\small\color{gris}\docpie\ · página \thepage}

\newcommand{\cabecera}[3]{%
  {\Huge\bfseries\color{acento}#1\par}\vspace{4pt}
  {\large #2\par}\vspace{2pt}
  {\color{gris}#3\par}\vspace{8pt}
  {\color{linea}\hrule height 0.6pt}\vspace{10pt}}
